\part{Lecture de fichiers externes}
\label{part:lecture}

\section{Introduction}
  Pour stocker les séries de mesure, AIDA utilise son propre format de 
fichiers de données, dit \textit{vecteur}. Les données issues de sources
diverses (en général des systèmes d'acquisition) doivent
donc être converties dans ce format de travail. Pour ce faire, il convient
d'utiliser le programme \texttt{mtv.py} (\textit{miscellaneous to vector})
(page : \pageref{lec-mtv})


\section{Lecture des fichiers d'acquisition : \texttt{mtv.py}}
\label{lec-mtv}
\texttt{mtv.py} permet de convertir des fichiers de données de différents
formats en fichiers-vecteur.

\subsection {Lancement}
\texttt{mtv.py} peut être lancé sous \texttt{aito} (onglet \texttt{Importation})
ou en ligne de commande dans un terminal. \\
Lancer la commande \fbox{ \texttt{mtv.py -h} } pour obtenir de l'aide.

\subsection {Manuel d'utilisation}
\fbox{ \texttt{mtv.py [\textit{options}] <\textit{source}>} }

\medskip
avec :
\begin{itemize}
\item[•] \textit{source} : argument obligatoire. \\
         \textit{source} est :
    \begin{itemize}
      \item \textbf {un fichier} : \\
            toute l'acquition est contenue dans ce seul fichier.
      \item \textbf {liste de fichiers} : \\
            l'acquition est découpée en plusieurs fichiers. La listes des
            fichiers doit être fournie dans l'ordre chronologique de
            l'acquisition.
      \item \textbf {un répertoire} : \\
            le répertoire ne contient que les fichiers des différentes
            parties de l'acquition. \\
            Les fichiers du répertoire sont lus dans l'ordre
            alphabétique : la chronologie de l'acquisition doit donc
            suivre l'ordre alphabétique des fichiers du répertoire.
      \item \textbf {une liste de répertoires} : \\
            la liste des répertoires est lue dans l'ordre fourni. Chaque
            répertoire doit suivre la logique décrite ci-dessus.
    \end{itemize}

\medskip
\item[•] \textit{options} :
    \begin{itemize}
      \item[] \texttt{-a/{}-{}-acqid <acqid>}           : acqid.
      \item[] \texttt{-p/{}-{}-plateforme <plateforme>} : nom long de la plateforme.
      \item[] \texttt{-c/{}-{}-campagne <campagne>}     : nom de la campagne (utile pour renseigner le fichier .inf).
      \item[] \texttt{-f/{}-{}-file <fichier>}          : fichier de configuration (chemin absolu ou relatif).
      \item[] \texttt{-t/{}-{}-type <type>}             : type de fichier à lire (csv\_like, ...).
      \item[] \texttt{-s/{}-{}-start <début>}           : AAAA[MMJJHHMMSS.sss] : date/heure de début des données à lire (inclus).
      \item[] \texttt{-e/{}-{}-end <fin>}               : AAAA[MMJJHHMMSS.sss] : date/heure de fin des données à lire (exclus).
      \item[] \texttt{-n/{}-{}-dry-run}                 : dry run (lecture sans écriture des vecteurs).
      \item[] \texttt{-F/{}-{}-force}                   : écrase les fichiers vecteurs existants.
      \item[] \texttt{-g/{}-{}-graphic}                 : trace les données.
      \item[] \texttt{-C/{}-{}-config <type>}           : aide sur la configuration associée au type de fichier (csv\_like, ...)
      \item[] \texttt{-k/{}-{}-keep}                    : conserve les valeurs hors \textit{[normal\_min, normal\_max]}.
      \item[] \texttt{-L/{}-{}-log\_dir <log\_dir>}     : répertoire des logs (fichier mtv.log).
      \item[] \texttt{-l/{}-{}-logger <logger>}         : fichier de log (chemin complet) dans lequel mtv rajoute ses logs.
      \item[] \texttt{-nc/{}-{}-no-cal}                 : pas de suffixe i\texttt{cal}.
      \item[] \texttt{-nk/{}-{}-no-checksum}            : pas de checksum.
      \item[] \texttt{-r/{}-{}-readers}                 : liste des readers.
      \item[] \texttt{-h/{}-{}-help}                    : affiche l'aide.
    \end{itemize}
\end{itemize}

\medskip
Remarques :

\begin{itemize}
\item la présence d'une des 2 options \texttt{acqid} ou \texttt{plateforme} est obligatoire.
\item si la plateforme est de type acquisition (\texttt{duree\_vecteur} dans le fichier \\
      \texttt{\$AIDA\_HOME/profils/<PROFIL>/plate-formes)} la présence de l'acqid est obligatoire.
\end{itemize}

\subsection {Format du fichier configuration :}
Le fichier de configuration permet à \texttt{mtv.py} de décoder le fichiers
d'acquisition qui peut être des différentes type. Ex :
\begin{itemize}
  \item fichier ASCII avec un en-tête, et des données rangées en colonnes séparées par un ou plusieurs séparateurs).
  \item fichier de type netCDF.
\end{itemize}

\medskip
Liste non exhaustive de fichiers de type ASCII : AROME/ARPEGE, CampbelRAY, G24D1M (Picarro), XSENS, GILL, LI7500, ...

\medskip
Le format du fichier configuration est de type "INI" : il est composé de paragraphes entre crochets suivi d'une liste
d'attributs de type «~champ = valeur~»

\medskip
De manière générale, la colonne est désignée par son numéro \textbf{commençant à 1} ou le nom de colonne définie dans la
ligne de nom de colonne si présente dans l'entête. Cette valeur est désignée :
\ col\_{\textless}\textit{1{\textgreater}} à \ col\textit{\_{\textless}N{\textgreater}} par la suite.

\medskip
Les paragraphes \ [col\_defaut] et [instru\_{\textless}N{\textgreater}] sont optionnels. Ils permettent de définir des
valeurs par défaut des variables. Si les valeurs ne sont pas définies dans le paragraphe de la variable alors les
valeurs des paragraphe [col\_defaut] et [instru\_N] sont utilisées. Cela évite de répéter les informations.

Si on veut utiliser pour un paramètre les valeurs des champs définies dans le paragraphe [instru\_N] alors ajouter un
champ instru = instru\_{\textless}N{\textgreater} dans la description du paramètre.

\medskip
Le paragraphe [msg\_{\textless}ID{\textgreater}] permet de décrire des fichiers où dans un même fichier on trouve des
lignes avec un nombre colonnes et/ou une fréquence des données différentes et dépendant d'une valeur de colonne. 4M
produit couramment ce genre de fichier. Pour les autres types de fichiers les champs définis dans
[msg\_{\textless}ID{\textgreater}] \ doivent être définis dans le paragraphe [structure]

\medskip
{\large \color{blue} [structure] (OBLIGATOIRE)}

\begin{center}
\tablefirsthead{}
\tablehead{\hline champ &
valeur &
Statut &
commentaire\\}
\tabletail{\hline}
\tablelasttail{}
\begin{supertabular}{|p{3.592cm}|p{2.116cm}|p{4.099cm}|p{5.393cm}|}
\hline
champ &
valeur &
Statut &
commentaire\\\hline
nb\_lignes\_entete &
{\itshape nombre} &
Optionnel

défaut = 0 &
Nombre de ligne de l'en{}-tête. \\\hline
numero\_ligne\_nom &
{\itshape nombre ${\geq}$ 1} &
Optionnel

défaut = None &
Numéro de la ligne commençant à 1 contenant les noms des variables correspondant aux colonnes de données. Le séparateur
des noms doit être identique au séparateur de colonnes de données.

Ce système devra être étendu pour traiter des lignes complexes du type : \textit{\textcolor{black}{\#variables : param1,
param2, .{\dots}}}\\\hline
numero\_ligne\_unit &
{\itshape nombre ${\geq}$ 1} &
Optionnel

défaut = None &
Numéro de la ligne contenant les unités des variables correspondant aux colonnes de données. Le séparateur des noms doit
être identique au séparateur de colonnes de données.\\\hline
\textbf{numero\_ligne}

\textbf{\_misval}  

(non implémenté) &
{\itshape nombre ${\geq}$ 1} &
Optionnel

défaut = None &
Numéro de la ligne contenant la valeur de données manquante des variables correspondant aux colonnes de données. Le
séparateur des noms doit être identique au séparateur de colonnes de données.\\\hline
numero\_ligne

\_borne\_inf  &
{\itshape nombre ${\geq}$ 1} &
Optionnel

défaut = None &
Numéro de la ligne contenant les valeurs inférieures des intervalles représentant les bornes des donées 2D. Le
séparateur des bornes doit être identique au séparateur de colonnes de données.\\\hline
numero\_ligne

\_borne\_sup  &
{\itshape nombre ${\geq}$ 1} &
Optionnel

défaut = None &
Numéro de la ligne contenant les valeurs supérieures des intervalles représentant les bornes des donées 2D. Le
séparateur des bornes doit être identique au séparateur de colonnes de données.\\\hline
{\color{black} separateurs\_colonnes} &
{\color{black} \textit{c}\textit{aractère}\textit{\textcolor[rgb]{0.4,0.4,1.0}{ }}\textit{ou \ ESPACE ou TAB}} &
{\color{black} OBLIGATOIRE}

{\color{black} Pas de valeur par défaut.} &
{\color{black} Liste des séparateurs existants dans le fichier de données.}


~

{\itshape\color{black} exemple : ESPACE , ; TAB}

~
\\\hline
{\color{black} separateurs\_multiples}

{\color{black} \_ignores} &
{\color{black} \textit{c}\textit{aractère}\textit{\textcolor[rgb]{0.4,0.4,1.0}{ }}\textit{ou \ ESPACE ou TAB}} &
{\color{black} Optionnel}

{\color{black} défaut = None} &
{\color{black} Séparateurs à ignorer s'ils se succèdent sur la ligne.}

~

{\itshape\color{black} Par exemple : si~ separateurs\_multiples\_ignores= TAB alors plusieurs TAB successifs sur la
ligne seront considérés comme un seul séparateur. }\\\hline
{\bfseries\color{black} separateur\_decimal} 

(non implémenté)&
{\itshape\color{black} caractère} &
{\color{black} Optionnel}

{\color{black} defaut=~'.' (point)} &
~
\\\hline
debut\_ligne\_a\_ignorer &
{\itshape chaîne} &
Optionnel

défaut=None &
Définition d{}'une chaîne de caractères de début de ligne à ignorer par la lecture.\\\hline
colonne\_msg\_id &
{\itshape col\_{\textless}N{\textgreater}}

{\itshape avec N ${\geq}$ 1}

{\itshape ou }

{\itshape col\_{\textless}NOM{\textgreater}}

~

~
 &
Optionnel.

défaut=None

~

Si colonne\_msq\_id n'est pas défini alors on ignore les paragraphes [msg\_{\textless}ID{\textgreater}]

et 

liste\_msg\_id &
~

Numéro de colonne contenant l'identifiant de message pour le cas des fichiers contenant des messages différents
(fréquence et nombre de colonnes).

Si ce champ est renseigné, renseigner les paragraphes [msg\_{\textless}ID1{\textgreater}],
[msg\_{\textless}ID2{\textgreater}], {\dots}.\\\hline
liste\_msg\_id &
\textit{{\textless}ID1{\textgreater}, {\textless}ID2{\textgreater},{\textless}ID3{\textgreater}}

{\itshape (chaîne de caractères)} &
Optionnel mais fortement recommandé dans le cas si colonne\_msg\_id est défini 

défaut = None &
Liste des msg\_ID possibles pour faire du contrôle.\\\hline
traite\_doublons  &
True/False &
Optionnel

défaut = True &
Si True : élimine les données en double (temps et valeur identiques)

Sinon toutes les données sont conservées.\\\hline
ignore\_doublons &
col\_5,col\_8 &
Optionnel

défaut = None &
Liste des colonnes ignorées par la recherche des doublons (séparateur = `,')\\\hline
\end{supertabular}
\end{center}

\medskip
{\itshape\color{black}
Exemple~de fichiers avec plusieurs séparateurs:}

{\itshape\color{black}
18/09/2016 12:00:01.23 CAX1 3.1, 5.2, 7.3}

\medskip
{\itshape\color{black}
separateurs\_colonnes = ESPACE ,}

{\itshape\color{black}
separateurs\_multiples\_ignores= ESPACE}

\medskip
{\texttt{\textbf{\large \color{blue} {[msg\_{\textless}ID{\textgreater}]}}} (existe si colonne\_msg\_id est défini)

\medskip
Si aucun [msg\_{\textless}ID{\textgreater}] n'est défini, les champs ci-dessous \textbf{doivent} se trouver dans le
paragraphe [structure]

\medskip
\begin{center}
\tablefirsthead{}
\tablehead{\hline champ &
valeur &
Statut &
commentaire\\}
\tabletail{\hline}
\tablelasttail{}
\begin{supertabular}{|p{3.504cm}|p{3.7cm}|p{3.642cm}|p{4.448cm}|}
\hline
Champ &
Valeur &
Statut &
Commentaire\\\hline
nb\_colonnes &
{\itshape \ nombre ${\geq}$ 1} &
Optionnel 

{\color{black} Pas utile si numero\_ligne\_nom est défini.}

~

{\color{black} Sinon \ il est OBLIGATOIRE}

~

{\color{black} défaut = None} &
Nombre de colonne de données \\\hline
{\bfseries\color{black} \textmd{temps} }

~

~
 &
{\itshape\color{black} source[d:f],...,source[d:f]

où source peut-être :

col\_{\textless}N{\textgreater}

NOMFIC

NOMREP} &
{\color{black} OBLIGATOIRE}

~
 &
{\color{black} NOMFIC et NOMREP sont des mots clefs.}

{\color{black} d et f représentent les indices de début de et fin de chaîne à décoder (cf. slice python).}

{\color{black} Le séparateur «~,~» est obligatoire et doit se retrouver dans format\_temps. }\\\hline
{\bfseries\color{black} \textmd{format\_temps} } &
{\itshape\color{black} fmt\_1,fmt\_2, ...}

~

{\itshape\color{black} Possibilité de définir plusieurs formats pour un même temps avec le séparateur | :}


{\itshape\color{black} fmt\_1|fmt\_2,fmt\_3,...}  &
{\color{black} OBLIGATOIRE} &
{\color{black} fmt\_1, fmt\_n sont des chaînes de caractères }

{\color{black} décrivant le codage du temps suivant strptime de Python enrichi des formats de délais : }

{\color{black} \%Ds (délai millisecondes),}

{\color{black} \%DS (délai secondes),}

{\color{black} \%DM (delai en minutes),}

{\color{black} \%DH(delai en heures),}

{\color{black} \%Dd(delai en jours),}

{\color{black} ainsi que du format temps GPS : }

{\color{black} \%G(nombre de semaines depuis le 06/01/1980),}


{\color{black} \%Di(secondes intercalaires à enlever)}\\\hline
{\color{black} UTC\_offset} &
{\itshape\color{black} +HHMM} &
{\color{black} Optionnel.}

{\color{black} défaut = +0000} &
{\color{black} UTC Offset (option \%z de strptime) }

~

{\color{black} meta-donnée qui se retrouve dans l'en{}-tête du fichier temps.}\\\hline
{\color{black} nom\_variable\_temps} &
{\itshape\color{black} chaîne } &
{\color{black} Optionnel}

~

{\color{black} défaut = rft\_time\_[nom\_fic config]\_[msg\_ID]} &
{\color{black} Nom du fichier vecteur temps si écriture vecteur , ou nom de variable si renvoi d'une structure de
données.}\\\hline
colonne\_checksum  &
{\itshape col\_{\textless}N{\textgreater} avec N ${\geq}$ 1}

{\itshape ou }

{\itshape col\_{\textless}NOM{\textgreater}} &
Optionnel

défaut = None &
Colonne contenant le check\_sum 

{\itshape ex : col\_7 , col\_CHECK}\\\hline
methode\_check\_sum &
{\itshape\color{black} nom de fonction} &
Optionnel

défaut = None &
Fonction connue de \texttt{mtv.py} avec range comme argument par ex :

{\itshape {\textless}xor\_checksum{\textgreater}:'{\textbackslash}x02', '{\textbackslash}x03'}\\\hline
freq\_ech &
{\itshape\color{black} nombre} &
Optionnel : par défaut la fréquence est calculée à partir des données. &
Fréquence du message en Hz.\\\hline
type\_ech &
{\itshape\color{black} chaîne} &
Optionnel~

défaut = None &
Type d{}'échantillonnage attendu du fichier :

regulier {\textbar} irregulier\\\hline
\end{supertabular}
\end{center}

\medskip
\texttt{\textbf{\large \color{blue} {[col\_defaut]}}} (optionnel)

\medskip
Si les champs ci-dessous ne sont pas définis dans le paragraphe décrivant la colonne alors on prend ces valeurs par
défaut.

\medskip
Les champs autorisés dans ce paragraphe sont :

\begin{itemize}
\item {\bfseries
normal\_min}
\item {\bfseries
normal\_max}
\item {\bfseries
misval\_in}
\item {\bfseries
misval\_out}
\item {\bfseries
invval\_in}
\item {\bfseries
invval\_out}
\item {\bfseries
format}
\item {\bfseries
unit}
\end{itemize}

\medskip
Pour le format de ces champs, voir le paragraphe \ (col\_{\textless}A{\textgreater}) de description des colonnes.

\medskip
\large      \color{blue}\textbf{[\texttt{instru\_N}]}
\normalsize \color{black} (optionnel)

\medskip
Ce paragraphe permet de définir des valeurs de champ liés à un instrument. La prise en compte de ces valeurs pour une
variable se fait via le champ instru= instru\_{\textless}N{\textgreater} dans le paragraphe colonne
(col\_{\textless}A{\textgreater})

\medskip
Les champs autorisés dans ce paragraphe sont :

\begin{itemize}
\item {\bfseries
position}
\item {\bfseries
chaîne acquisition}
\item {\bfseries
module\_{\textless}1{\textgreater} à module\_{\textless}n{\textgreater} }
\end{itemize}

\medskip
Pour le format de ces champs, voir le paragraphe \ (col\_{\textless}A{\textgreater}) de description des colonnes.

\medskip
\large      \color{blue}\textbf{[\texttt{col\_A:\{col\_B\{@msg\_msg\_id\}\}}]}
\normalsize \color{black} (Optionnel)

\medskip
Exemples de nom de paragraphe descriptif de variable :

\medskip
[col\_1] : paramètre de colonne 1

[col\_TPR] : paramètre correspondant à la colonne nommée TPR dans l'en-tête

[col\_2 : col\_10] : paramètre 2D entre les colonnes 2 et 10

[col\_10@msg\_1] : paramètre en colonne 10 pour le msg\_id = 1

\medskip
\begin{center}
\tablefirsthead{}
\tablehead{\hline champ &
valeur &
Statut &
commentaire\\}
\tabletail{\hline}
\tablelasttail{}
\begin{supertabular}{|p{2.9099998cm}|p{3.783cm}|p{4.058cm}|p{4.4500003cm}|}
\hline
{\bfseries\color{black} Champ} &
{\bfseries\color{black} Valeur} &
{\bfseries\color{black} Statut} &
{\bfseries\color{black} Commentaire}\\\hline
ignore  &
{\itshape True/ False} &
Optionnel

défaut = False &
Si 'True' \ la colonne de données correspondante est lue mais la variable n'est pas dans la structure de données
renvoyée (vecteur, ou autre...)\\\hline
{\color{black} controle } &
{\color{black} \textit{{\textless}nom\_fonction{\textgreater}:~{\textless}valeur{\textgreater}}} &
Optionnel

défaut = None &
Permet de spécifier un contrôle sur \ une colonne du fichier. 

Exemples des fonctions implémentées:

{\itshape {\textless}egal{\textgreater} :{}'L751'}\\\hline
instru &
{\itshape instru\_{\textless}N{\textgreater}} &
Optionnel

défaut = None

~

Vérifier que le paragraphe [instru\_{\textless}N{\textgreater}] \ existe. &
Correspond au paragraphe 

instru\_{\textless}N{\textgreater}\\\hline
nom\_variable &
{\itshape chaîne} &
Optionnel

défaut = col\_A ou nom de la colonne &
nom du fichier vecteur si écriture vecteur , ou nom de variable si renvoi d'une structure de données.\\\hline
name &
{\itshape chaîne} &
Optionnel.\newline
défaut = «~col\_A du fichier NOMFIC~» &
Nom long explicite de la variable\\\hline
english\_name &
{\itshape chaîne} &
Optionnel

défaut = «~col\_A from file NOMFIC~» &
Nom long explicite en anglais de la variable\\\hline
unit &
{\itshape chaîne} &
Optionnel mais fortement recommandé.

défaut = unknown &
Unité de la variable après application de la «~law~» éventuelle, \ acceptable par libudunits2\\\hline
format  &
{\itshape chaîne} &
Optionnel

défaut=ieeefloat32 &
Format connu de AIDA\\\hline
law &
{\itshape chaîne} &
Optionnel.

defaut = None &
Loi à appliquer selon norme AIDA 
(voir doc package readers,
 
module abstract\_reader,
 
rechercher 'law')

Si pas de law écriture directe d'un vecteur « cal »

Si présence de law écriture du vecteur BRUT « raw » et du vecteur « cal »

~

~
\\\hline
normal\_min &
{\itshape nombre} &
Optionnel

defaut=None &
Borne min de l'intervalle de validité de la variable (après application d'une éventuelle law).\\\hline
normal\_max &
{\itshape nombre} &
Optionnel

defaut=None &
Borne max de l'intervalle de validité de la variable (après application d'une éventuelle law).\\\hline
misval\_in &
{\itshape chaîne} &
Optionnel

défaut = NULL &
Valeur manquante dans le fichier d'entrée.

Si non NULL la présence d'un champ vide doit générer un signalement.\\\hline
misval\_out &
{\itshape chaîne} &
Optionnel

défaut = NaN &
Valeur manquante dans les données de sortie.\\\hline
invval\_in &
{\itshape chaîne} &
Optionnel

défaut = NULL &
Valeur invalide (différente de misval\_in) dans le fichier d'entrée\\\hline
invval\_out &
{\itshape chaîne} &
Optionnel

défaut = misval\_out &
Valeur invalide dans le tableau ou vecteur produit (champ invval du fichier vecteur)\\\hline
bornes  &
{\itshape (x{1} ,y{1} ); ...; (x{n },y{n} )}

(cf documentation de référence)

ou

{\itshape auto}&
Optionnel mais fortement recommandé si la variable est un tableau 1D

défaut = None &
Ce champ précise les bornes de chacun des éléments du tableau. S'il vaut auto, les valeurs des bornes sont automatiquement lues dans le fichier de données en utilisant les champs numero\_ligne\_borne\_inf et numero\_ligne\_borne\_sup.\\\hline
unit\_bornes &
{\itshape chaîne} &
Optionnel mais fortement recommandé

défaut = None &
Indique l'unité de l'attribut bornes\\\hline
position &
{\itshape chaîne} &
Optionnel mais fortement recommandé !

défaut = None &
Position de la mesure\\\hline
orientation &
{\itshape chaîne} &
Optionnel

défaut = None &
Lorsque le paramètre contenu dans un vecteur est une grandeur physique orientée

(un angle, une distance sur un axe, une composante d'un vecteur), ce champ contient une

description de cette orientation\\\hline
chaîne\_acquisition &
\textit{module\_1 {\dots}. module\_n} &
Optionnel mais fortement recommandé !

défaut = None &
Description de la chaîne d'acquisition selon norme AIDA (liste des modules)\\\hline
module\_{\textless}1{\textgreater} à module\_{\textless}n{\textgreater} &
{\itshape {\textless}type\_module{\textgreater}:{\textless}fabricant{\textgreater}:{\textless}Modèle{\textgreater}
:{\textless}Numéro de série{\textgreater}:{\textless}Numéro de référence}

{\itshape interne{\textgreater}:{\textless}date et lieu de calibrage ou étalonnage{\textgreater}} &
Optionnel mais fortement recommandé !

défaut = None &
Description d{}'un module d'acquisition selon norme AIDA\\\hline
information &
{\itshape chaîne } &
Optionnel

défaut = None &
Commentaire libre\\\hline
\end{supertabular}
\end{center}

\medskip
{\color{black}
Il est possible d'ajouter des champs à cette liste mais ils ne seront pas écrits dans les vecteurs. Ils seront dans la
structure de données renvoyée.}

\medskip
Les champs suivants sont présents dans les vecteurs produits mais déterminés automatiquement :

\medskip
\begin{itemize}
\item \textbf{dim\_array}~: est déterminé \ à partir de la dimension de chaque variable 
\item \textbf{obs\_min }et \textbf{obs\_max}~: sont calculées automatiquement
\item \textbf{freq\_ech} : calculée à partir du \ temps lu (inverse de la médiane de \texttt{${\Delta}$}t)
\item \textbf{famille}\textbf{~}: déterminée à partir des 3 premières lettres du champ \textbf{nom\_variable}
\item \textbf{origine~ }vaut «~lecture~»
\item \textbf{fct\_origine : }est la copie de la commande mtv complète~: «~ mtv {\dots}.~»
\end{itemize}

%%%\medskip
%%%\begin{itemize}
%%%\item[] 
%%%\bigskip
%%%\end{itemize}

\subsection{Evolution des fichiers de configuration dans le temps}
\label{sec:histo_config}
\texttt{mtv.py} ne prend pas en charge la gestion de l'évolution des fichiers
de configuration au cours du temps. Cela doit être fait en amont. AIDA dispose de
l'outil \texttt{histo\_config.py}. Cet outil utilise des fichiers .histo et
permet de connaître le fichier de configuration à utiliser par mtv en fonction de la
date.

\medskip
\subsubsection{Format des fichiers .histo}
Les fichiers .histo sont des fichiers de type .ini contenant des {\bfseries [sections]} et des lignes de type {\bfseries attribut = valeur} où :

{\bfseries section} = coupleur

{\bfseries attribut} = intervalle de temps => YYYY-mm-ddTHH:MM:SS~~YYYY-mm-ddTHH:MM:SS

{\bfseries valeur} = fichier de configuration utilisé par mtv

\subsubsection{Exemple de fichier .histo}
\# METEOPOLE-FLUX

\# Historique des fichiers configuration de la plateforme toulouse-flux

\#

\#

[campbel1]

2012-06-15T00:00:00 2015-06-20T23:59:59 = METEOPOLE\_FLUX.toulouse-flux.campbel1.V0.config

\# ajout du KT15

2015-06-21T00:00:00 2030-12-31T23:59:59 = METEOPOLE\_FLUX.toulouse-flux.campbel1.V1.config



[campbel2]

2012-06-15T00:00:00 2016-01-20T23:59:59 = METEOPOLE\_FLUX.toulouse-flux.campbel2.V0.config

\# ajout de la ventilation du CNR1

2016-01-21T00:00:00 2030-12-31T23:59:59 = METEOPOLE\_FLUX.toulouse-flux.campbel2.V1.config



[campbel3]

2012-06-15T00:00:00 2030-12-31T23:59:59 = METEOPOLE\_FLUX.toulouse-flux.campbel3.V0.config



\subsection{Gestion des «~log~»}
Le programme mtv.py écrit ses logs aux moyen du module python loggin.

Plusieurs loggers sont utilisés :

\begin{itemize}
\item mtv.log : ce fichier est systématiquement créé avec le niveau DEBUG. Par défaut, il se trouve sous /tmp mais son
emplacement peut être défini par l'option {}-L/-{}-log\_dir. Le fichier mtv.log est un logger de type
RotatingFileHandler~: dès que le fichier atteint une taille de 100Mo il est renommé et un autre fichier se crée. On
conserve ainsi une archive de 10 fichiers en mode tournant.
\item stdout : l'écriture sur stdout se fait en mode INFO par défaut. La variable d'environnement MTV\_LEVEL (DEBUG
{\textbar} INFO {\textbar} WARNING {\textbar} ERROR) permet de modifier le niveau de log
\item l'utilisateur peut également spécifier un fichier de log via l'option -l/-{}-logger dans lequel mtv ajoutera ses
logs en mode DEBUG.
\end{itemize}

\subsection[Architecture ]{Architecture} 
Le package «~readers~» contient les fonctions de «~lecture données~» et
«~contrôle temps~». Il est appelable par n'importe quel programme Python.
Il n'est pas spécifique ni de AIDA ni du format vecteur. Pour plus de
commodité nous installons néanmoins ce package dans AIDA.

\medskip
\texttt{mtv.py} est basé sur le package de lecture de données «~readers~».

\medskip
«~readers~» se trouve sous \texttt{\$AIDA\_HOME/src/python3/packages/readers}

\medskip
\texttt{mtv.py} se trouve sous \texttt{\$AIDA\_HOME/src/python3/lecture/mtv}

\clearpage
\subsection[Exemples de description du temps]{Exemples de description du temps dans le fichier configuration}
\tiny
\begin{flushleft}
\tablefirsthead{}
\tablehead{\normalsize\hline
{\bfseries Type de fichier} &
{\bfseries Format de la datation} &
{\bfseries Description dans fichier configuration }\\\tiny}
\tabletail{\hline}
\tablelasttail{}
\begin{supertabular}{|p{4.109cm}|p{5.3cm}|p{6.284cm}|}
\hline
{\bfseries \normalsize Type de fichier} &
{\bfseries \normalsize Format de la datation} &
{\bfseries \normalsize Description dans fichier configuration }\\\hline
AROME/ARPEGE &
2016050113 12.43 {\dots}.  &
separateur\_colonnes = ESPACE

temps = col\_1

format\_temps =\%Y\%m\%d\%H\\\hline
Fichier Campbel 4M (avant MEF)  &
17/05/2016 00:00:00.060 CMI1 {\textquotedbl}05/17/2016 00:00:00{\textquotedbl} , {\dots}.

~
 &
separateur\_colonnes = ,

temps = col\_1[0:23]

format\_temps =\%d/\%m/\%Y \%H:\%M:\%S.\textcolor{black}{\%f}\textcolor{red}{ }

(respecter l'espace entre date et heure dans \ 

format\_temps)\\\hline
Fichier GILLV 4M  &
\ 31/05/2016 00:30:00.022 GIV1 220 9 56 17049  &
separateur\_colonnes = ESPACE

temps = col\_1, col\_2

format\_temps=\%d/\%m/\%Y,

\%H:\%M:\%S.\textcolor{black}{\%f }\\\hline
Fichier GILL et XSENS nacelle turbulente &
Nom de fichier : tr\_09080842.GIL ou tr\_09080842.XSE

Dans les fichiers :

~

08:08:42.261,3, {\dots}. &
On n'a pas la date complète ni dans le nom du fichier, ni dans les données. Passer l'Acqid en argument de \texttt{mtv.py}

~

Dans le fichier configuration on aura :

separateur\_colonnes = ,

temps = col\_1

format\_temps = \%H:\%M:\%S.\textcolor{black}{\%f}\textcolor{red}{ }

~

{\color{black} \textbf{ou} si on a la date complète dans nom de fichier (Eric peut nous faire ça ..)}

{\color{black} comme par exemple : \ tr\_2016\_09080842.GIL}

~

{\color{black} separateur\_colonnes = ,}

{\color{black} temps = NOMFIC[3:12], col\_1}

{\color{black} format\_temps =\%Y\_\%m\%d, \%H:\%M:\%S.\%f }

~
\\\hline
Essais KH20~(4M): &
Dans les fichiers :

{\textquotedbl}2015-07-02 12:59:14.9{\textquotedbl},382355,{\dots}..

{\textquotedbl}2015-07-02 12:59:15{\textquotedbl},382356, {\dots}.. &
separateur\_colonnes = ,

temps = col\_1

format\_temps =~'{}'\%Y-\%m{}-\%d\%H:\%M:\%S.\%f'{}'

~

Attention parfois on n'a pas les décimales ...\\\hline
RADOME vidage Degreanne &
20141029000600 70 2.4 &
separateur\_colonnes = ESPACE

temps = col\_1

format\_temps =\%Y\%m\%d\%H\%M\%S\\\hline
PICARRO  &
DATE \ \ \ \ \ \ \ \ \ \ \ \ TIME \ \ \ \ \ \ \ \ FRAC\_DAYS\_SINCE\_JAN1

2015-11-05 \ \ \ \ 13:04:00.000 \ 308.54444444 &
separateur\_colonnes = \ ESPACE

temps = col\_DATE,col\_TIME

format\_temps =\%Y-\%M{}-\%d,\%H :\%M :\%S\\\hline
Message AIPOV de l'AIRINS &
\$AIPOV,133306.6012,...

\$AIPOV,133306.6112,... &
separateur\_colonnes = ,

temps = col\_2

format\_temps =\%H\%M\%S.\%f\\\hline
Températures de brillance issues de rad2Tb  &
2013-05-30,07:54:01.392, 7.90039,284.696...

2013-05-30,07:54:01.552, 7.90043,284.714...

2013-05-30,07:54:01.712, ... &
separateur\_colonnes = ,

temps = col\_Date,col\_ClockTime

format\_temps =\%Y-\%m-\%d,\%H :\%M :\%S.f \\\hline
PACS &
1,\#2014-03-26 13:31:06\#,633,\#2014-03-26 13:31:06\#,633,1000 &
separateur\_colonnes = ,

temps = col\_2

format\_temps= \#\%Y-\%m-\%d \%H :\%M :\%S\#\\\hline
Licor7500 &
DATAH\ \ Time\ \ Date\ \ Sequence Number 

DATA\ \ 04:39:30:549\ \ 2012-09-26\ \ 298870 &
separateur\_colonnes = TAB

temps = col\_Date,col\_Time

format\_temps =\%Y-\%m-\%d,\%H :\%M :\%S

~
\\\hline
RS format \ 

RS\_FLEDT\_G4253301\_

20121031085852.tsv

~

~
 &
time Pscl T RH v u ...

0000 28372.5 291.85 81 1.84228 3.32355 2 ...

0002 28368.4 291.666 85.172 1.85628 2.9116 ... &
separateur\_colonnes = ESPACE

temps = NOMFIC~[18:32], col\_time

format\_temps =\%Y\%m\%d\%H\%M\%S,\%Ds\\\hline
\end{supertabular}
\end{flushleft}
\normalsize

\subsection{Fichier d'acquisition au format netCDF}
\fbox{ \texttt{mtv.py -C netcdf} } \\
permet d'obtenir de l'aide sur la structure du fichier de configuration attendu
par le reader \texttt{netcdf}.
